\section{Pillar 4 --- Fault Isolation and Blast Radius in Distributed Systems}
\textit{(Theoretical basis for \sq{2})}

A foundational principle of distributed systems design is that failures should be contained
within the smallest possible boundary \citep{tran2022}. The concept of blast radius --- the
scope of system impact from a single failure event --- is central to reliability
architecture in microservices and cloud-native systems.

In a single-\vm{} shared execution environment, the blast radius of any individual job
failure is potentially the entire system. A job that exhausts memory can trigger
out-of-memory process termination affecting all other jobs on the same host. In a
Kubernetes pod model, each job runs in an isolated namespace with enforced resource limits,
meaning a failing pod is terminated by the \textsc{oom} killer without directly affecting
other pods on the same node.

This theoretical blast radius reduction is one of the primary architectural motivations for
the Kubernetes migration. Whether this reduction holds in practice at the workflow execution
level --- and by how much it reduces execution time variance and improves job success rate
for concurrently running pipelines --- is an empirical question that existing literature
at the infrastructure level does not answer for application-level workflow systems.

\textbf{Research connection:} \sq{2} empirically tests the blast radius reduction at the
workflow execution level, measuring whether pod isolation contains failures and reduces
execution time variance under equivalent concurrent workload conditions.
